\section{Finite Product Transfer and Event Loss}

The setup in Section 3 defines \(\ptheta(\Lang\mid x)\) for any finite score
table over \(\Y^T\).  The product-transfer computation below requires more
structure.  Assume the score admits a finite-context local factorization.  Let \(C\) be a
finite context state space with initial context \(c_0\), deterministic update
\(U:C \times \Y \to C\), and local potentials
\(\phi_t(x,c_{t-1},y_t)\) such that
\[
  S_{\theta}(x,y) = \sum_{t=1}^T \phi_t(x,c_{t-1},y_t).
\]
Let \(G_t(c,j)=\exp\phi_t(x,c,j)\).  Let a complete DFA for \(\Lang\) be
\((Q,q_{\mathrm{start}},F,\delta)\).  Define the product transfer
\[
M_t^{\Lang}[(c,q),(c',q')]
= \sum_{j \in \Y} G_t(c,j)
  \mathbf{1}[c'=U(c,j)]\mathbf{1}[q'=\delta(q,j)] .
\]
Let \(u_0\) be one at \((c_0,q_{\mathrm{start}})\) and let
\(b_{\Lang}(c,q)=\mathbf{1}[q \in F]\).

The finite-context assumption covers standard first-order linear-chain CRFs and
fixed higher-order Markov contexts.  An arbitrary finite score table can be
represented with a sufficiently large context or table construction, but that
fallback is exponential in general and does not support an efficiency or
scaling claim.  Product-state scaling statements therefore apply only to the
chosen finite-context implementation.

\begin{theorem}
Under the finite setup above, including finite-context local factorization and
a complete deterministic DFA,
\[
  u_0 M_1^{\Lang}\cdots M_T^{\Lang} b_{\Lang}
  = \Z_{\theta,\Lang}(x).
\]
\end{theorem}

\begin{proof}
Fix a label sequence \(y=(y_1,\ldots,y_T)\).  The deterministic context update
and complete DFA transition induce a unique trajectory
\[
  (c_0,q_0),(c_1,q_1),\ldots,(c_T,q_T),
  \qquad q_0=q_{\mathrm{start}},
\]
where \(c_t=U(c_{t-1},y_t)\) and \(q_t=\delta(q_{t-1},y_t)\).  Along this path,
the product of local transfer weights is
\[
  \prod_{t=1}^T G_t(c_{t-1},y_t)
  = \exp\left(\sum_{t=1}^T \phi_t(x,c_{t-1},y_t)\right)
  = \exp S_\theta(x,y).
\]
Conversely, any nonzero term in the expansion of
\(u_0M_1^\Lang\cdots M_T^\Lang b_\Lang\) selects labels that satisfy the same
deterministic context and DFA updates.  If multiple labels induce the same
product transition, the transfer entry sums their weights, and distributivity
of the finite matrix product still assigns each label sequence exactly its own
product of local weights.  The terminal vector \(b_\Lang\) keeps exactly the
trajectories with \(q_T\in F\), which by DFA correctness are exactly
\(\Lang_T\).  Thus the product transfer sums \(\exp S_\theta(x,y)\) over
\(y\in\Lang_T\), which is \(\Z_{\theta,\Lang}(x)\).
\end{proof}

\begin{proposition}
If scores are finite and differentiable and \(\Z_{\theta,\Lang}(x)>0\), then
\[
\nabla_{\theta}[-\log \ptheta(\Lang \mid x)]
= \E_{\ptheta(y\mid x)}[\nabla_{\theta}S_{\theta}(x,y)]
 - \E_{\ptheta(y\mid x,y\in\Lang_T)}[\nabla_{\theta}S_{\theta}(x,y)].
\]
\end{proposition}

\begin{proof}
Because \(\Y^T\) is finite, differentiation commutes with summation.  Therefore
\[
 \nabla_\theta \log \Z_\theta(x)
 = \sum_{y\in\Y^T}\ptheta(y\mid x)\nabla_\theta S_\theta(x,y)
\]
and, when \(\Z_{\theta,\Lang}(x)>0\),
\[
 \nabla_\theta \log \Z_{\theta,\Lang}(x)
 = \sum_{y\in\Lang_T}\ptheta(y\mid x,y\in\Lang_T)
   \nabla_\theta S_\theta(x,y).
\]
Since \(-\log\ptheta(\Lang\mid x)=\log\Z_\theta(x)-\log\Z_{\theta,\Lang}(x)\),
subtracting the two identities gives the result.
\end{proof}

This proposition explains the event loss as a computable training signal.  It
does not imply that task accuracy, F1, calibration, or constrained-decoding
quality improves.

\paragraph{Diagnostic interpretation.}
The scalar \(\ptheta(\Lang\mid x)\) is attached to the original posterior, not
to the selected decoded sequence.  It can therefore be evaluated even when a
constrained decoder returns a legal output.  Low event mass should be
interpreted as posterior inconsistency with a specified rule; whether that
inconsistency predicts task error is an empirical diagnostic question.
